\begin{table}[!t]
\centering
\caption{Experimental platforms and common evaluation settings.}
\label{tab:experimental_setup}
\footnotesize
\renewcommand{\arraystretch}{1.10}
\setlength{\tabcolsep}{4pt}

\begin{tabularx}{\columnwidth}{p{0.30\columnwidth}X}
\toprule
\textbf{Item} & \textbf{Configuration} \\
\midrule

\rowcolor{gray!10}
\multicolumn{2}{l}{LLM configuration} \\
LLM &
DeepSeek API with \texttt{deepseek-v4-pro}; temperature $=0$. \\

\midrule
\rowcolor{gray!10}
\multicolumn{2}{l}{Processor and software environment} \\
CPU candidates &
Ariane/CVA6 (RV64), Ibex (RV32), and Leon3 (SPARC V8). \\

Compilation &
RISC-V GCC for Ariane/CVA6 and Ibex; \texttt{sparc-elf} GCC for Leon3.
CPU-only and CPU--accelerator implementations use identical compiler
settings for each application. \\

Embedded runtime &
Benchmark-specific accelerator-control and runtime code is manually
implemented using the ESP programming model and is held fixed during
evaluation. \\

\midrule
\rowcolor{gray!10}
\multicolumn{2}{l}{Accelerator generation and estimation} \\
HLS &
Vitis/Vivado HLS 2021.2; all five primary DSE runs use a 13.333-ns
target clock (75~MHz). \\

CPU estimation &
Sniper-based CPU performance estimation. \\

Area estimation &
A common analytical SoC-area model is used for hardware/software
partitioning and DSE under identical technology assumptions. \\

\midrule
\rowcolor{gray!10}
\multicolumn{2}{l}{System-level exploration and validation} \\
SoC generation &
ESP-based CPU--accelerator SoC configuration and integration. \\

DSE &
HyperMapper-based exploration over processor, cache, and accelerator
configuration parameters under application-specific area constraints. \\

FPGA validation &
Xilinx Zynq UltraScale+ \texttt{xczu9eg-ffvb1156-2-e};
ESP runtime clock of 75~MHz. \\

\bottomrule
\end{tabularx}
\end{table}